% Source: https://artofproblemsolving.com/community/c3h2361784_how_to_find_the_next_term_of_any_sequence
% Stable identity: aops-topic-2361784-q01
The first ten terms of a sequence are $8,3,4,9,3,9,8,3,4,3$.
\begin{enumerate}
\item[(a)] Assuming that the $n$th term is given by a polynomial in $n$ of degree at most $9$, find the eleventh term.
\item[(b)] A student claims that the method ``take differences until one number is left, assume that it stays constant, and work back'' always finds the next term of any sequence. In what sense is this true, and in what sense is it wrong?
\end{enumerate}
